\section{Explicit periodic geometry and a genuine inducing set}
Two elementary cycles have zero winding. The normal nearest-neighbor orbit has two flights, total length
\begin{equation}\label{eq:two-three}
 L_2=2(1-2R),\qquad
 L_3=3(1-\sqrt3 R)
\end{equation}
is the length of the inner equilateral three-cycle on centers $0,a,b$. Its contact angles are $\pi/6,5\pi/6,3\pi/2$ and its outgoing $p$ coordinates are all $-1/2$. The reflection cosines are respectively $1$ and $\sqrt3/2$. The equilateral segments lie inside their elementary triangle and outside the three vertex disks; any other lattice center is at distance at least $\sqrt3/2$ from that triangle. Hence all other disks are strictly avoided.

\begin{lemma}[A winding four-cycle]\label{lem:winding}
For every $R\in I$ there is an analytic angle $\theta_R\in(0.92,0.967)$ satisfying
\begin{equation}\label{eq:theta}
 f_R(\theta)=2R\sin\theta-\sin\theta\cos\theta
                       +\frac{\sqrt3}{2}\cos(2\theta)=0.
\end{equation}
The itinerary with centers $0,a,a-b,2a-b,a$ and normals at angles
\[
 -\theta_R,\quad\pi+\theta_R,\quad\theta_R,\quad\pi-\theta_R
\]
is a physical four-collision periodic orbit modulo $\Lambda$, with winding $a$. All its incidence cosines exceed $1/2$ and its clearance from every third disk exceeds $1/200$. Rotation through $\pi/3$ gives another cycle with winding $b$ and the same length.
\end{lemma}
\begin{proof}
Put $c=\cos\theta$, $s=\sin\theta$. Its first four contact points and the translated fifth point are
\begin{align*}
 q_0&=(Rc,-Rs),&q_1&=(1-Rc,-Rs),\\
 q_2&=(1/2+Rc,-\sqrt3/2+Rs),&q_3&=(3/2-Rc,-\sqrt3/2+Rs),\\
 q_4&=q_0+a.
\end{align*}
Two segments are horizontal. Write $d=2Rc-1/2$, $h=\sqrt3/2-2Rs$. The other two are $(d,-h)$ and $(d,h)$. The scalar equation $d\sin2\theta+h\cos2\theta=0$ is precisely \eqref{eq:theta}. Since $d,h>0$, it makes the diagonal direction the reflection of the horizontal direction at the intervening normal. The remaining reflection equations follow by symmetry, with cosine $c$ at every contact. The total length is
\begin{equation}\label{eq:fourlength}
 L_4=2(1-2Rc)+2\sqrt{d^2+h^2}.
\end{equation}

Here are reproducible rational bounds for the inequalities and absence of occlusion. Divide $I$ into 32 closed intervals of width $1/1600$. On each, bisect the endpoint equations \eqref{eq:theta} between $9/10$ and $1$ to width at most $10^{-10}$. The interval derivative
\[
 2R\cos\theta-\cos2\theta-\sqrt3\sin2\theta
\]
is less than $-1/2$. The endpoint signs and positivity of $\partial_R f$ show that the two endpoint root brackets enclose the root for the entire parameter interval. The resulting angles lie in $(0.92,0.967)$; $d,h>0$ and $c>1/2$ on every interval.

For clarity, the arithmetic of this finite certificate uses no floating root solver. Square roots are enclosed by integer square roots at denominator $10^{22}$. For sine and cosine use the Taylor polynomials through degrees 25 and 24 at the rational midpoint, with errors at most $|x|^{26}/26!$ and $|x|^{25}/25!$, respectively; add the interval radius, since both functions are 1-Lipschitz. All interval operations are rational.

For each segment check centers $ia+jb$, $|i|,|j|\le4$, excluding its two endpoints. The distance from the rational midpoint segment is computed by clamping the rational orthogonal projection to $[0,1]$. Subtract the endpoint and center enclosure radii and the largest obstacle radius. Each lower bound exceeds $1/200$. Every contact point has norm less than 3; a center outside the index box has norm at least $5\sqrt3/2$, so it too has distance greater than the obstacle radius from every segment. The complete deterministic calculation is specified in \texttt{tools/certify\_winding.py}; it obtains lower bounds greater than $0.00925$ for clearance and $0.569$ for incidence, and a derivative upper bound below $-0.731$. The signs and strict margins prove the assertions on the full intervals, not only at sampled radii. The implicit function theorem gives analytic dependence and uniqueness in the printed root interval.
\end{proof}

Choose $\epsilon=1/200$. In $(\alpha,p)$ coordinates let $Y_R$ be the union of the open squares of half-side $\epsilon$ centered at
\begin{equation}\label{eq:Y}
 (0,0),\quad(\pi/6,-1/2),\quad
 (-\theta_R,\sin\theta_R),\quad
 (\pi/3-\theta_R,\sin\theta_R).
\end{equation}
Angles are taken modulo $2\pi$. The squares are disjoint and
$\nu(Y_R)=4\epsilon^2/\pi$. Let $r_R$ be the first return time to $Y_R$ and $F_R=T_R^{r_R}$ on its recurrent full-measure domain, with invariant probability $\nu_R^Y=\nu|_{Y_R}/\nu(Y_R)$. Define
\begin{equation}\label{eq:induced}
 \kappa_R=\sum_{j<r_R}\kappa_{\rm coll}\circ T_R^j,
 \qquad \tau_R^Y=\sum_{j<r_R}\tau_R\circ T_R^j.
\end{equation}
Poincar\'e recurrence and invariance give this induced probability; they do not give an exponential return tail, a Markov partition, or a differentiable family of spectral projectors.

\begin{proposition}[Full augmented lattice at zero roof frequency]\label{prop:lattice}
The four cycles above give fixed points of $F_R$ with data $(\kappa_1,\kappa_2,r,\text{return period})$ equal to the rows of
\begin{equation}\label{eq:matrix}
 V=\begin{pmatrix}0&0&2&1\\0&0&3&1\\1&0&4&1\\0&1&4&1\end{pmatrix},
                  \qquad |\det V|=1.
\end{equation}
Consequently a circle-valued coboundary identity
$e^{i(u\cdot\kappa+s r)}=e^{ic}g/(g\circ F_R)$ that can be evaluated at these fixed points forces $u=s=c=0$ modulo $2\pi$.
\end{proposition}
\begin{proof}
Only one contact of each cycle belongs to \eqref{eq:Y}. For the four-cycle the successive $p$ coordinates are $s,s,-s,-s$. The same is true after rotation. The two positive-$p$ unselected contacts are separated from the selected centers; the smallest relevant angular separation is $2\pi/3-2\theta_R>0.16$. All negative-$p$ contacts are separated in $p$ from the selected winding contacts, and also from $0,-1/2$. The remaining contacts of the two- and three-cycles are separated by their displayed angles. Thus the physical first return counts are $2,3,4,4$, not freely assigned labels. The four selected contacts are interior to their squares, so the same finite return branches exist in neighborhoods. Direct subtraction of the first two rows gives determinant of magnitude one. Evaluation first gives $2s=c=3s$, hence $s=c=0$; the last two rows then give $u_1=u_2=0$.
\end{proof}

The qualification about evaluation is essential. A merely measurable transfer function is not automatically defined at a periodic point. A spectral application requires its own regularity or Liv\v{s}ic argument. Moreover, adding a roof twist $b\tau_R^Y$ to this proposition leaves $bL_j$ in every periodic equation. The four rows do not eliminate it. Joint roof nonarithmeticity and returned temporal nonintegrability remain distinct requirements. As one smaller unconditional conclusion, \eqref{eq:two-three} excludes a continuous cohomology $\tau_R=c+g-g\circ T_R$: the two periodic averages are $1-2R$ and $1-\sqrt3R$, which differ by $(2-\sqrt3)R>0$.

